% ============================================================================
% 第二章：相关工作
% ============================================================================

\section{运动想象 BCI 系统}

\begin{frame}{运动想象的神经生理基础}
    \begin{block}{事件相关去同步化 (ERD/ERS)}
        \begin{itemize}
            \item 运动想象时大脑感觉运动皮层产生 ERD 现象
            \item $\mu$频段 (8--13 Hz) 和$\beta$频段 (13--30 Hz) 功率显著降低
            \item 空间分布具有任务特异性
        \end{itemize}
    \end{block}
    
    \vspace{0.5cm}
    \begin{alertblock}{关键挑战}
        ERD 的时间动态特性存在显著的\textbf{个体间差异}，使得固定时间窗策略难以适应所有被试
    \end{alertblock}
\end{frame}

\begin{frame}{BCI 系统的基本架构}
    \begin{center}
    \begin{tikzpicture}[node distance=0.8cm]
        \node[draw, rectangle, rounded corners, fill=xmu!10, minimum width=2.5cm, minimum height=0.8cm] (data) {数据采集};
        \node[draw, rectangle, rounded corners, fill=xmu!10, minimum width=2.5cm, minimum height=0.8cm, right=of data] (preprocess) {预处理};
        \node[draw, rectangle, rounded corners, fill=xmu!20, minimum width=2.5cm, minimum height=0.8cm, right=of preprocess] (feature) {特征提取};
        \node[draw, rectangle, rounded corners, fill=xmu!30, minimum width=2.5cm, minimum height=0.8cm, right=of feature] (select) {时间窗选择};
        \node[draw, rectangle, rounded corners, fill=xmu!10, minimum width=2.5cm, minimum height=0.8cm, right=of select] (classify) {分类};
        
        \draw[->, thick] (data) -- (preprocess);
        \draw[->, thick] (preprocess) -- (feature);
        \draw[->, thick] (feature) -- (select);
        \draw[->, thick] (select) -- (classify);
    \end{tikzpicture}
    \end{center}
    
    \vspace{0.5cm}
    \begin{exampleblock}{本研究定位}
        时间窗选择属于\textbf{特征选择}子任务，直接影响特征提取的质量和分类性能
    \end{exampleblock}
\end{frame}

% ============================================================================
% EEG 特征提取与分类
% ============================================================================
\section{EEG 特征提取与分类}

\begin{frame}{共空间模式 (CSP)}
    \begin{block}{核心思想}
        寻找最优空间滤波器，使得两类任务的方差差异最大化
    \end{block}
    
    \vspace{0.5cm}
    \begin{itemize}
        \item \textbf{协方差矩阵}: $\mathbf{C}_c = \frac{1}{N_c} \sum_{i \in \text{class}_c} \mathbf{X}_i \mathbf{X}_i^\top$
        \item \textbf{广义特征值}: $\mathbf{C}_1 \mathbf{w} = \lambda (\mathbf{C}_1 + \mathbf{C}_2) \mathbf{w}$
        \item \textbf{特征提取}: $\mathbf{f} = \log \left( \text{var}(\mathbf{W} \mathbf{X}) \right)$
    \end{itemize}
    
    \vspace{0.5cm}
    \begin{alertblock}{特点}
        \begin{itemize}
            \item 计算高效，二分类任务表现优异
            \item 对噪声敏感，仅适用于二分类
        \end{itemize}
    \end{alertblock}
\end{frame}

\begin{frame}{支持向量机 (SVM)}
    \begin{block}{最大间隔分类}
        \[
        \min_{\mathbf{w}, b} \frac{1}{2} \|\mathbf{w}\|^2 + C \sum_{i=1}^N \xi_i
        \]
    \end{block}
    
    \vspace{0.5cm}
    \begin{itemize}
        \item 通过寻找最大间隔超平面实现分类
        \item 小样本场景下表现稳健
        \item 正则化参数$C$控制过拟合风险
    \end{itemize}
    
    \vspace{0.5cm}
    \begin{exampleblock}{CSP-SVM 流水线}
        CSP 特征提取 + 线性 SVM 分类 $\rightarrow$ 本研究的基础分类器
    \end{exampleblock}
\end{frame}

% ============================================================================
% 时间窗优化方法
% ============================================================================
\section{时间窗优化方法}

\begin{frame}{固定时间窗方法}
    \begin{table}[h]
        \centering
        \footnotesize
        \begin{tabular}{l|l|l}
            \hline
            \textbf{策略} & \textbf{方法} & \textbf{局限} \\
            \hline
            刺激后固定延迟 & 0.5--3.5 秒 & 忽略个体差异 \\
            \hline
            基于 ERD 峰值 & 选择 ERD 最强时段 & 需额外校准 \\
            \hline
            经验时间窗 & 2--4 秒 & 缺乏理论依据 \\
            \hline
        \end{tabular}
    \end{table}
    
    \vspace{0.5cm}
    \begin{alertblock}{局限性}
        \begin{itemize}
            \item 无法适应个体间/个体内差异
            \item 需要手动调参，依赖研究者经验
        \end{itemize}
    \end{alertblock}
\end{frame}

\begin{frame}{基于优化的时间窗选择}
    \begin{enumerate}
        \item \textbf{贝叶斯优化}
        \begin{itemize}
            \item 构建高斯过程代理模型
            \item 计算复杂度$O(n^3)$，难以处理高维空间
        \end{itemize}
        
        \item \textbf{网格搜索/随机搜索}
        \begin{itemize}
            \item 实现简单，计算效率低
        \end{itemize}
        
        \item \textbf{遗传算法}
        \begin{itemize}
            \item 收敛速度慢，需设计编码和变异策略
        \end{itemize}
        
        \item \textbf{稀疏组空间模式}
        \begin{itemize}
            \item 联合优化时间窗与 CSP
            \item 增加计算复杂度，正则化参数需调优
        \end{itemize}
    \end{enumerate}
\end{frame}

\begin{frame}{时间窗优化的评估标准}
    \begin{itemize}
        \item \textbf{分类准确率}: 最直接的评估指标
        \item \textbf{Cohen's Kappa}: 考虑随机猜测的影响
        \[
        \kappa = \frac{p_o - p_e}{1 - p_e}
        \]
        \item \textbf{计算效率}: 搜索时间 + 推理延迟
        \item \textbf{稳定性}: 跨被试/跨 session 的标准差
    \end{itemize}
\end{frame}

% ============================================================================
% 强化学习在 BCI 中的应用
% ============================================================================
\section{强化学习在 BCI 中的应用}

\begin{frame}{强化学习基础}
    \begin{block}{马尔可夫决策过程 (MDP)}
        五元组 $(\mathcal{S}, \mathcal{A}, \mathcal{R}, \mathcal{P}, \gamma)$
    \end{block}
    
    \vspace{0.5cm}
    \begin{itemize}
        \item \textbf{Q-Learning}: 维护状态 - 动作价值表 (Q 表)
        \item \textbf{Q 值更新}: $Q(s, a) \leftarrow Q(s, a) + \alpha [r + \gamma \max_{a'} Q(s', a') - Q(s, a)]$
        \item \textbf{收敛性保证}: 有限状态空间下以概率 1 收敛到最优策略
    \end{itemize}
    
    \vspace{0.5cm}
    \begin{exampleblock}{优势}
        无需环境模型、实现简单、低维状态空间下高效
    \end{exampleblock}
\end{frame}

\begin{frame}{深度 Q 网络 (DQN)}
    \begin{block}{核心思想}
        使用深度神经网络近似 Q 函数
    \end{block}
    
    \vspace{0.5cm}
    \begin{itemize}
        \item \textbf{经验回放}: 打破样本相关性
        \item \textbf{目标网络}: 提升训练稳定性
        \item \textbf{改进方法}: Double Q-Learning、Dueling 架构等
    \end{itemize}
    
    \vspace{0.5cm}
    \begin{alertblock}{局限性}
        \begin{itemize}
            \item 模型复杂（百万级参数）
            \item 需要大量训练数据
            \item 可解释性差（黑盒）
        \end{itemize}
    \end{alertblock}
\end{frame}

\begin{frame}{强化学习在 BCI 中的应用}
    \begin{enumerate}
        \item \textbf{动态窗口选择} (Chen et al., 2019)
        \begin{itemize}
            \item SSVEP 识别，Q-Learning 学习最优策略
        \end{itemize}
        
        \item \textbf{多智能体特征选择} (MARS, Zhang et al., 2021)
        \begin{itemize}
            \item 联合优化空间、频谱和时间特征
            \item 计算复杂度高，需要大量数据
        \end{itemize}
        
        \item \textbf{自适应 BCI 控制}
        \begin{itemize}
            \item 调整分类器参数、选择最优 BCI 范式
        \end{itemize}
    \end{enumerate}
\end{frame}

% ============================================================================
% 本章小结
% ============================================================================
\section{本章小结}

\begin{frame}{本章小结}
    \begin{block}{文献综述发现}
        \begin{itemize}
            \item 现有强化学习方法（如 MARS）计算复杂度高
            \item 多智能体设置和深度网络结构带来巨大开销
            \item 难以应用于资源受限的在线 BCI 系统
        \end{itemize}
    \end{block}
    
    \vspace{0.5cm}
    \begin{alertblock}{核心研究问题}
        在运动想象时间窗优化这一\textbf{状态空间有限}（136 个状态）的特定任务中，是否可以采用更简洁的\textbf{表格型 Q-Learning}，在避免复杂性的同时，保留强化学习的自适应能力？
    \end{alertblock}
\end{frame}
